% =====================================================================
% An End-to-End Multi-Objective Reinforcement Learning Framework for
% UAV Path Planning with Dynamic Obstacle Avoidance and Energy Efficiency
% IEEE iskeleti — Overleaf'e kopyalayip kullanabilirsiniz.
% =====================================================================
\documentclass[journal]{IEEEtran}
\usepackage[utf8]{inputenc}
\usepackage{amsmath,amssymb,graphicx,booktabs,cite,url,algorithm,algorithmic}
\graphicspath{{figures/}}

\begin{document}

\title{An End-to-End Multi-Objective Reinforcement Learning Framework for
UAV Path Planning with Dynamic Obstacle Avoidance and Energy Efficiency}

\author{Ad~SOYAD,~\IEEEmembership{}% TODO
\thanks{A. Soyad, Bölüm, Üniversite, Şehir, Türkiye (e-posta: ...).}}

\maketitle

\begin{abstract}
% TODO: paper/makale_taslagi.md'deki özeti sonuçlarla doldurup çevirin.
\end{abstract}

\begin{IEEEkeywords}
UAV, path planning, multi-objective reinforcement learning, dynamic obstacle
avoidance, energy efficiency, Pareto optimization.
\end{IEEEkeywords}

\section{Introduction}
% Motivasyon, boşluk, katkılar (bkz. taslak §1)

\section{Related Work}
\subsection{Classical Path Planning}
\subsection{RL-Based UAV Navigation}
\subsection{Multi-Objective Reinforcement Learning}
\subsection{UAV Energy Models}

\section{Problem Formulation}
We model the task as a multi-objective MDP
$\mathcal{M}=(\mathcal{S},\mathcal{A},P,\mathbf{r},\gamma)$ with
$\mathbf{r}(s,a)=[r_{\text{goal}},r_{\text{energy}},r_{\text{safety}}]^\top$.

\subsection{UAV Kinematics}
\begin{equation}
\mathbf{v}_{t+1}=(1-\lambda)\mathbf{v}_t+\mathbf{a}_t\Delta t,\qquad
\|\mathbf{v}\|\le v_{\max},\ \|\mathbf{a}\|_\infty\le a_{\max}.
\end{equation}

\subsection{Propulsion Energy Model}
Following Zeng \emph{et al.}~\cite{zeng2019energy}:
\begin{equation}
\begin{aligned}
P(V)={}&P_0\!\left(1+\frac{3V^2}{U_{\mathrm{tip}}^2}\right)
+P_i\!\left(\sqrt{1+\frac{V^4}{4v_0^4}}-\frac{V^2}{2v_0^2}\right)^{\!1/2}\\
&+\tfrac12 d_0\rho s A V^3 + mg[V_z]^+ .
\end{aligned}
\end{equation}

\subsection{Reward Vector and Scalarization}
\begin{equation}
R_t=\mathbf{w}^\top\mathbf{r}_t,\qquad \mathbf{w}\sim\mathrm{Dirichlet}(\mathbf{1})
\ \text{(training)}.
\end{equation}

\section{Proposed Framework}
\subsection{Preference-Conditioned Policy}
\subsection{Learning Algorithms}
\subsection{Classical Baselines}

\section{Experimental Setup}
% Arena, engeller, metrikler, hiperparametre tabloları

\section{Results and Discussion}
% Tablo/şekiller: results/figures ve results/metrics'ten
\begin{figure}[t]\centering
\includegraphics[width=\columnwidth]{pareto_front.png}
\caption{Approximate Pareto front obtained from a single preference-conditioned policy.}
\end{figure}

\section{Conclusion}

\bibliographystyle{IEEEtran}
\begin{thebibliography}{9}
\bibitem{zeng2019energy} Y. Zeng, J. Xu, and R. Zhang, ``Energy minimization for wireless communication with rotary-wing UAV,'' \emph{IEEE Trans. Wireless Commun.}, vol.~18, no.~4, pp. 2329--2345, 2019.
\bibitem{schulman2017ppo} J. Schulman \emph{et al.}, ``Proximal policy optimization algorithms,'' arXiv:1707.06347, 2017.
\bibitem{haarnoja2018sac} T. Haarnoja \emph{et al.}, ``Soft actor-critic,'' in \emph{Proc. ICML}, 2018.
\bibitem{khatib1986apf} O. Khatib, ``Real-time obstacle avoidance for manipulators and mobile robots,'' \emph{Int. J. Robot. Res.}, 1986.
\bibitem{raffin2021sb3} A. Raffin \emph{et al.}, ``Stable-Baselines3: Reliable reinforcement learning implementations,'' \emph{JMLR}, 2021.
\bibitem{yang2019envelope} R. Yang, X. Sun, and K. Narasimhan, ``A generalized algorithm for multi-objective reinforcement learning and policy adaptation,'' in \emph{Proc. NeurIPS}, 2019.
\end{thebibliography}

\end{document}
